
Table~\ref{tab:specA_set2_nva_main} reports the baseline estimation
results for the aggregate capital specification, focusing on the
effective wage interaction (Set 2) on the NVA basis. Full results
for all 32 specifications---spanning both interaction sets, both
value-added bases, and both structural break modes---are reported
in Appendix~\ref{app:specA_full_results}.

\paragraph{Cointegration verdict.}
The cointegration test triad confirms the presence of a stable
cointegrating relationship across all 32 estimated models. The
Engle-Granger residual ADF test rejects the null hypothesis of no
cointegration in every specification ($t_{\text{ADF}} \leq -2.32$,
with augmented lag selection via AIC). The Shin
\citeyearpar{Shin1994} LM statistic ($C_{\text{Shin}} \leq 0.06$)
and the Vogelsang-Wagner Type 2 statistic fail to reject the null
hypothesis of cointegration across all specifications. This mutual
confirmation establishes that the cointegrating relationship between
$lnY^{NFC}$ and aggregate private gross capital stocks is robust to the
choice of wage share measure, interaction specification, and
structural break treatment.

\paragraph{Capital elasticity.}
The capital elasticity $\hat{\theta}_0$ is positive and highly
significant across all specifications, ranging from 0.55 to 0.61
with standard errors of 0.02--0.08. The point estimates are
remarkably stable across wage share measures, interaction sets, and
break modes, indicating that the long-run output-capital elasticity
is robustly identified. The inclusion of Bai-Perron step dummies
reduces $\hat{\theta}_0$ by approximately 0.02 in most
specifications, reflecting the absorption of level shifts into the
deterministic component.

\paragraph{Failure of the standard wage interaction (Set 1).}
In the standard interaction specifications (Set 1), the wage
interaction coefficient $\hat{\theta}_1$ fails to achieve
statistical significance under the custom Monte Carlo fixed-$b$
critical values across all eight wage share measures and both break
modes. The point estimates range from $-0.19$ to $+0.06$ with
standard errors of 0.10--0.14, yielding $|t|$-statistics below the
90th percentile Monte Carlo threshold of 2.96. This failure
demonstrates that the raw, unconditioned wage share---whether
computed over GVA or NVA, and whether using the baseline NFC measure
or the Shaikh-Tonak supervisory-adjusted measures---does not capture
the accumulation-channel dynamics predicted by the theoretical
framework.

\paragraph{Success of the effective wage interaction (Set 2).}
When the wage interaction is scaled by productive accumulation
$(1 - \iota_{t-1})$, where $\iota_{t-1}$ is the lagged share of
intellectual property products in total capital, the effective wage
interaction coefficient $\hat{\theta}_1$ becomes statistically
significant with the theoretically predicted negative sign in the
Shaikh-Tonak Tier 2 NVA specifications. In model M2.8a (no
structural break dummies), $\hat{\theta}_1 = -0.21$ ($t = -2.63$,
significant at the 10\% level under Monte Carlo critical values);
in model M2.8b (with Bai-Perron step dummies), $\hat{\theta}_1 =
-0.35$ ($t = -5.00$, significant at the 5\% level). The negative
sign implies that higher effective wage pressure---wage pressure
net of the unproductive intellectual property component---reduces
the marginal capacity yield of capital. Economically, high effective
wage pressure forces defensive mechanization, increasing the organic
composition of capital and depressing the transformation elasticity.

The contrast between Set 1 and Set 2 is sharp: the standard wage
share fails to mediate the capital-output relationship, while the
effective wage share---purged of the unproductive IPP
component---captures the distributive channel through which class
struggle over newly created value shapes the pace of capacity
accumulation. This result is consistent with the theoretical
prediction that only wage pressure over net value added, net of
unproductive components, enters the firms' cost-minimization
calculus and thereby mediates the transformation of capital into
productive capacity.

\paragraph{Structural break robustness.}
The inclusion of Bai-Perron step dummies does not alter the
qualitative conclusions. In Set 2 NVA specifications, the effective
wage interaction remains significant with the same sign after
controlling for endogenous structural breaks. The break years
identified by the Bai-Perron procedure are reported in
Appendix~\ref{app:specA_full_results}. The stability of
$\hat{\theta}_1$ across break modes confirms that the distributive
mediation result is not an artifact of omitted structural breaks.

%%%% Results Table %%%%%
\begin{table}[h!]
\centering
\caption{Specification A, Set 2: Effective Wage Interaction (NVA Basis)}
\label{tab:specA_set2_nva_main}
\small
\begin{tabular}{p{4.2cm}p{1.2cm}p{1.2cm}p{1.2cm}p{1.2cm}p{1.2cm}p{1.2cm}p{1.2cm}p{1.2cm}}
\toprule
 & \multicolumn{2}{c}{NFC NVA Baseline} & \multicolumn{2}{c}{Clean Corp Tier 1} & \multicolumn{2}{c}{Clean Corp Tier 2} & \multicolumn{2}{c}{ST Tier 2 NVA} \\
\cmidrule(lr){2-3} \cmidrule(lr){4-5} \cmidrule(lr){6-7} \cmidrule(lr){8-9}
 & (1) & (2) & (3) & (4) & (5) & (6) & (7) & (8) \\
\midrule
$\hat{\theta}_0$ (Capital Elasticity) & 0.58$^{***}$ & 0.56$^{***}$ & 0.58$^{***}$ & 0.65$^{***}$ & 0.57 & 0.56 & 0.57 & 0.55 \\
 & (0.02) & (0.06) & (0.02) & (0.06) & (0.02) & (0.08) & (0.02) & (0.06) \\
$\hat{\theta}_1$ (Effective Wage Pressure) & 0.00 & -0.17 & -0.04 & -0.30$^{**}$ & -0.19 & -0.49 & -0.21 & -0.35 \\
 & (0.10) & (0.07) & (0.08) & (0.06) & (0.08) & (0.13) & (0.08) & (0.07) \\
\midrule
Step Shifts ($D_{T_b}^{\text{step}}$) & No & Yes & No & Yes & No & Yes & No & Yes \\
$C_{\text{Shin}}$ [$H_0$: coint.] & 0.06 & 0.04 & 0.07$^{*}$ & 0.03 & 0.06 & 0.04 & 0.06 & 0.04 \\
$S_{\text{Type 2}}$ [$H_0$: coint.] & 10.76 & 9.07 & 11.05 & 9.06 & 7.50 & 10.89 & 6.13 & 13.18 \\
$t_{\text{ADF}}$ [$H_0$: no coint.] & -2.32 & -3.06 & -2.22 & -4.07$^{**}$ & -1.84 & -2.58 & -2.00 & -2.72 \\
Bandwidth ratio $b$ & 0.298 & 0.169 & 0.280 & 0.165 & 0.272 & 0.237 & 0.259 & 0.237 \\
Sample Size $T$ & 94 & 94 & 94 & 94 & 76 & 76 & 76 & 76 \\
\bottomrule
\end{tabular}
\begin{minipage}{\textwidth}
\scriptsize
\textit{Notes:} Dependent variable: $y_t^{\text{NFC}}$ (log NFC real GVA). Estimated via Cointegrating Multivariate Polynomial Regression (CMPR) IM-OLS \citep{VogelsangWagner2024}. Step Shifts "No": no structural break dummies; "Yes": Bai-Perron step dummies included. All interaction terms are FWL-orthogonalized against base capital regressors and step dummies. The effective wage interaction is $\omega_{t-1}(1-\iota_{t-1}) \times k_{\text{cap},t}$, where $\iota_{t-1}$ is the lagged IPP share of capital. Significance: $^*p<0.10$, $^{**}p<0.05$, $^{***}p<0.01$. Stars for coefficient rows and cointegration tests are derived from a 100,000-pass Monte Carlo simulation suite mapping the exact DGP variant and fixed-$b$ Wald tables. Standard errors in parentheses.
\end{minipage}
\end{table}




\subsubsection{Baseline Transformation Elasticity and Capacity Utilization Estimates}

\subsection{United States: Heterogeneous Capital Stock Estimates}

\subsubsection{Heterogeneous Capital Stocks Empirical Estimates}
\label{sec:431:specB_results}

Specification B disaggregates the capital stock into non-residential
structures ($k^{\text{NRC}}_t$) and the mechanization gap
($\kappa_t = k^{\text{ME}}_t - k^{\text{NRC}}_t$), replacing the
single aggregate capital regressor of Specification A with two
distinct capital components. The cointegrating regression takes the
form
\begin{equation}
	y_t = \alpha + \theta_{\text{NRC}}\, k^{\text{NRC}}_t
	+ \theta_{\kappa}\, \kappa_t
	+ \theta_{\omega\kappa}\, \xi^{\omega}_t
	+ \theta_{\kappa p}\, \xi^{p}_t
	+ u_t,
	\label{eq:431:specB_model}
\end{equation}
where $\xi^{\omega}_t$ is the Frisch-Waugh-Lovell (FWL)
orthogonalized wage-capital interaction and $\xi^{p}_t$ is the FWL
orthogonalized relative-price interaction. The FWL orthogonalization
partials out the base capital regressors and any detected structural
break step dummies before estimating the interaction coefficients,
as described in \cref{sec:421:estimation_strategy}. Two interaction
specifications are estimated. Set 1 uses the standard wage
interaction $\xi^{\omega}_t = \kappa_t \times \omega_{t-1}$; Set 2
uses the effective wage interaction $\xi^{\omega}_t = \kappa_t
\times \omega_{t-1}(1 - \iota_{t-1})$, where $\iota_{t-1}$ is the
lagged share of intellectual property products in total capital. All
models are estimated via the integrated modified OLS (IM-OLS)
estimator \citep{VogelsangWagner2014}, whose partial-sum
transformation and augmentation procedure for cointegrating
polynomial regressions are derived in
\cref{sec:421:estimation_strategy}. The estimation sample spans
1931--2024 ($T = 94$ for standard wage measures; $T = 77$ for
Shaikh-Tonak Tier 2 measures, which require the supervisory
compensation ratio available from 1948).

\paragraph{Theoretical motivation for the NVA basis.}
The wage share entering the interaction term is computed over either
gross value added (GVA) or net value added (NVA). Under the
classical-Marxian decomposition, $\text{GVA} = \text{CFC} + v + s$
and $\text{NVA} = v + s$, where CFC denotes consumption of fixed
capital. Because CFC represents pre-existing constant capital
transferred to output rather than newly created value, it lies
outside the zero-sum distributive conflict over surplus. During
intensive mechanization, the ratio $\text{CFC}/\text{GVA}$ rises
mechanically, depressing GVA-basis wage shares even when the rate of
surplus-value extraction ($s/v$) is unchanged. The NVA basis
isolates the distributive balance from this mechanization
distortion. Additionally, the Shaikh-Tonak Tier 2 measures strip
supervisory overhead from compensation, isolating productive
variable capital \citep{ShaikhTonak1994}.

\paragraph{Cointegration verdict.}
The cointegration test triad---comprising the Engle-Granger
residual ADF test, the Shin \citeyearpar{Shin1994} LM test, and
the Vogelsang-Wagner Type 2 statistic---is applied across all 24
estimated models. The construction of these test statistics and
their fixed-$b$ critical values are described in
\cref{sec:421:estimation_strategy}. The Engle-Granger residual ADF
test rejects the null hypothesis of no cointegration in all
specifications ($t_{\text{ADF}} \leq -4.16$, $p < 0.05$). The Shin
LM test ($C_{\text{Shin}} \leq 0.05$) and the Vogelsang-Wagner
Type 2 statistic ($S_{\text{Type 2}} \leq 16.25$) fail to reject
the null hypothesis of cointegration across all specifications. This
mutual confirmation establishes that the cointegrating relationship
between output and disaggregated capital is stable, and that the
Bai-Perron step dummies absorb historical regime shifts without
introducing spurious trend contamination.

\paragraph{Set 1 results: standard wage interaction.}
Tables~\ref{tab:app_specB_set1_gva} and
\ref{tab:app_specB_set1_nva} in Appendix~\ref{app:E} report the
Set 1 estimates. Three findings emerge. First, the structures
coefficient $\hat{\theta}_{\text{NRC}}$ is positive and highly
significant across all specifications ($0.35$--$0.62$,
$p < 0.01$), confirming that non-residential structures constitute
the primary capacity-yield channel. Second, the mechanization gap
coefficient $\hat{\theta}_{\kappa}$ is positive but statistically
insignificant in most GVA-basis specifications, becoming significant
only in the Tier 2 NVA specification with break dummies
($\hat{\theta}_{\kappa} = 1.00$, $p < 0.05$). Third, and most
critically, the standard wage interaction
$\hat{\theta}_{\omega\kappa}$ is statistically insignificant in all
Set 1 specifications. The point estimates range from $-13.07$ to
$+4.01$ with standard errors of $3.10$--$7.41$, yielding
$|t|$-statistics below 2.0 in all cases. The unadjusted wage share,
whether computed over GVA or NVA, fails to capture the
accumulation-channel dynamics predicted by the theory.

\paragraph{Set 2 results: effective wage interaction.}
Tables~\ref{tab:app_specB_set2_gva} and
\ref{tab:specB_set2_nva} report the Set 2 estimates, in which the
wage interaction is scaled by productive accumulation
$(1 - \iota_{t-1})$. The effective wage interaction
$\hat{\theta}_{\omega\kappa}$ becomes statistically significant
with the theoretically predicted negative sign. In the
Shaikh-Tonak Tier 1 GVA specification without break dummies
(M2.17a), $\hat{\theta}_{\omega\kappa} = -15.99$ ($t = -3.24$,
$p < 0.01$); with break dummies (M2.18b),
$\hat{\theta}_{\omega\kappa} = -12.29$ ($t = -2.91$,
$p < 0.05$). In the Tier 2 GVA specification (M2.21a), the
coefficient reaches $-18.45$ ($t = -4.78$, $p < 0.01$). The
NVA-basis results in Table~\ref{tab:specB_set2_nva} confirm the
pattern: $\hat{\theta}_{\omega\kappa} = -12.77$ ($t = -3.35$,
$p < 0.01$) in M2.15a and $-15.43$ ($t = -3.98$, $p < 0.01$) in
M2.23a.

The negative sign of $\hat{\theta}_{\omega\kappa}$ implies that
higher effective wage pressure reduces the marginal capacity yield
of mechanization. Economically, when productive labor commands a
larger share of newly created value, firms respond with defensive
mechanization---substituting machinery for labor to restore cost
competitiveness---but this substitution raises the organic
composition of capital without a proportional expansion of
productive capacity. The mechanization gap $\kappa_t$ widens, but
the transformation elasticity $\theta(\omega_{t-1}) =
\theta_{\text{NRC}} + \theta_{\kappa} +
\theta_{\omega\kappa}\,\omega_{t-1}(1-\iota_{t-1})$ declines,
depressing the capacity yield per unit of capital.

\paragraph{Relative price interaction.}
The relative-price interaction $\hat{\theta}_{\kappa p}$ is
statistically insignificant in all specifications, with point
estimates ranging from $-1.53$ to $+1.38$ and no consistent sign
pattern. This suggests that relative capital price movements do not
systematically mediate the mechanization-capacity channel in the
U.S. data over the full sample period. The result is consistent with
the theoretical prediction that relative price effects operate
primarily through the investment allocation channel (which determines
$\kappa_t$) rather than through the capacity yield channel (which
determines $\theta$).

\paragraph{Robustness across break specifications.}
Comparing specifications without break dummies against variants
controlling for Bai-Perron break step dummies reveals that the
structural break step dummies absorb level shifts without materially
altering the slope coefficients. The structures coefficient
$\hat{\theta}_{\text{NRC}}$ typically declines by $0.03$--$0.13$
when break dummies are included, reflecting the absorption of level
shifts into the deterministic component. The effective wage
interaction coefficient retains its sign and significance in all
Set 2 specifications, confirming that the distributive mediation
result is not an artifact of omitted structural breaks.

%%%% Results Table %%%%%
\begin{table}[h!]
	\centering
	\caption{Specification B, Set 2: Effective Wage Interaction (NVA Basis)}
	\label{tab:specB_set2_nva_main}
	\small
	\begin{tabular}{p{3.0cm}p{1.25cm}p{1.25cm}p{1.25cm}p{1.25cm}p{1.5cm}p{1.5cm}p{1.5cm}p{1.5cm}}
		\toprule
		& \multicolumn{2}{c}{NFC NVA Baseline} & \multicolumn{2}{c}{Clean Corp Tier 1} & \multicolumn{2}{c}{Clean Corp Tier 2} & \multicolumn{2}{c}{ST Tier 2 NVA} \\
		\cmidrule(lr){2-3} \cmidrule(lr){4-5} \cmidrule(lr){6-7} \cmidrule(lr){8-9}
		& (1) & (2) & (3) & (4) & (5) & (6) & (7) & (8) \\
		\midrule
		$\hat{\theta}_{\text{NRC}}$  (Structures $k^{\text{NRC}}$) & 0.61$^{***}$ & 0.58$^{***}$ & 0.60$^{***}$ & 0.60$^{***}$ & 0.54$^{***}$ & 0.54$^{***}$ & 0.54$^{***}$ & 0.51$^{***}$ \\
		& (0.02) & (0.08) & (0.02) & (0.08) & (0.01) & (0.01) & (0.01) & (0.01) \\
		$\hat{\theta}_{\kappa}$ (Mechanization Gap $\kappa$) & 0.32$^{*}$ & 0.39 & 0.32$^{*}$ & 0.39 & 0.31$^{***}$ & 0.31$^{***}$ & 0.33$^{***}$ & 0.06 \\
		& (0.13) & (0.29) & (0.13) & (0.29) & (0.06) & (0.06) & (0.06) & (0.09) \\
		$\hat{\theta}_{\omega\kappa}$ (Wage Interaction) & -12.77$^{**}$ & -10.13$^{**}$ & -11.08$^{**}$ & -7.93$^{**}$ & -14.97$^{***}$ & -14.97$^{***}$ & -17.05$^{***}$ & -16.76$^{***}$ \\
		& (3.81) & (3.43) & (3.48) & (2.96) & (2.01) & (2.01) & (2.34) & (2.55) \\
		$\hat{\theta}_{\kappa p}$ (Price Interaction) & 0.18 & 0.92 & 0.69 & 1.41$^{*}$ & -0.31 & -0.31 & -0.79$^{**}$ & -0.71$^{**}$ \\
		& (0.51) & (0.69) & (0.63) & (0.75) & (0.33) & (0.33) & (0.32) & (0.31) \\
		\midrule
		Step Shifts ($D_{T_b}^{\text{step}}$) & No & Yes & No & Yes & No & Yes & No & Yes \\
		$C_{\text{Shin}}$ [$H_0$: coint.] & 0.03 & 0.02 & 0.03 & 0.03 & 0.02 & 0.02 & 0.02 & 0.02 \\
		$S_{\text{Type 2}}$ [$H_0$: coint.] & 14.97 & 5.93 & 15.37 & 6.37 & 8.96 & 8.96 & 10.28 & 8.70 \\
		$t_{\text{ADF}}$ [$H_0$: no coint.] & -4.33$^{*}$ & -4.58$^{**}$ & -4.34$^{*}$ & -4.60$^{**}$ & -4.24$^{*}$ & -4.24$^{*}$ & -4.71$^{**}$ & -4.57$^{**}$ \\
		Bandwidth ratio $b$ & 0.185 & 0.173 & 0.193 & 0.181 & 0.208 & 0.208 & 0.208 & 0.208 \\
		Sample Size $T$ & 94 & 94 & 94 & 94 & 76 & 76 & 76 & 76 \\
		\bottomrule
	\end{tabular}
	\begin{minipage}{\textwidth}
		\scriptsize
		\textit{Notes:} Dependent variable: $y_t^{\text{NFC}}$ (log NFC real GVA). Estimated via Cointegrating Multivariate Polynomial Regression (CMPR) IM-OLS \citep{VogelsangWagner2024}. Step Shifts "No": no structural break dummies; "Yes": Bai-Perron step dummies included. All interaction terms are FWL-orthogonalized against base capital regressors ($k^{\text{NRC}}_t, \kappa_t$) and step dummies. The effective wage interaction is $\xi^\omega_t = \omega_{t-1}(1-\iota_{t-1}) \times \kappa_t$, where $\kappa_t = k^{\text{ME}}_t - k^{\text{NRC}}_t$ is the capital composition ratio (mechanization gap) and $\iota_{t-1}$ is the lagged IPP share of capital. Significance: $^*p<0.10$, $^{**}p<0.05$, $^{***}p<0.01$. Stars for coefficient rows and cointegration tests are derived from a 100,000-pass Monte Carlo simulation suite mapping the exact DGP variant and fixed-$b$ Wald tables. Standard errors in parentheses.
	\end{minipage}
\end{table}